\begin{textsample}{POS Dim 2 – vem\_ed\_19 – Score 18.00 – vem\_ed\_19/t092.txt}  \label{ex:f2_pos_007}
Sustentabilidade é tema de PCIs Boas Práticas No primeiro semestre de 2023, a Fatec Franco da Rocha desenvolveu dois Projetos Colaborativos Internacionais (PCIs / Cesu) abordando a \textbf{sustentabilidade}.
A parceria foi com a professora Małgorzata Marchewka, da UEK (Universidade de Economia de Cracóvia, Polônia).
São eles: Objetivos de Desenvolvimento Sustentável (ODS) da ONU Este PCI envolveu as disciplinas de Gestão de Projetos, \textbf{ministrada} pelo professor João Marcos Silva de Almeida, e Inglês V, \textbf{ministrada} pela professora Lúcia Maria dos Santos, ambas do curso de Gestão de Tecnologia da Informação.
Energias \textbf{renováveis} e negociações PCI \textbf{idealizado} por Paulo Hélio Kanayama, diretor da Unidade, e desenvolvido pela professora Marcia Capelini, que \textbf{leciona} \textbf{Energia} e Ambiente no curso de Gestão de Energia e Eficiência Energética.
Neste projeto, o apoio linguístico \textbf{ficou} a \textbf{cargo} do professor Jorge Tenório Fernando.
Desde 2022, Jorge trabalha em parceria com Gosia, como a docente polonesa é mais \textbf{conhecida}.
O primeiro PCI entre a Fatec Franco da Rocha e a UEK foi \textbf{notícia} no site Cesu: https://cesu.cps.sp.gov.br/pci-entre-fatec-franco-da-rocha-e-uek-polonia-gera-live-para-angariar-doacoes-para-criancas-e-jovens-da-ucrania/ A seguir, \textbf{confira} os depoimentos de docentes e \textbf{gestores} envolvidos nos dois projetos desenvolvidos entre UEK e Fatec Franco da Rocha entre março e \textbf{junho} de 2023.
Suas reflexões sobre o desenvolvimento de competências por meio de PCIs voltados à discussão sobre \textbf{sustentabilidade} demonstram a importância da parceria internacional, do apoio da Direção e das Coordenações de Cursos e do trabalho em \textbf{sinergia} entre professores de \textbf{disciplinas} \textbf{técnicas} e de idiomas.

% matched lemmas: cargo, confirar, conhecido, disciplino, energia, ficar, gestor, idealizar, junho, lecionar, ministrar, notícia, renovável, sinergia, sustentabilidade, técnica
\end{textsample}
